\chapter{Introdução}
\label{chap:introducao}

O avanço da Inteligência Artificial (IA) e, em particular, das técnicas de visão computacional tem possibilitado o desenvolvimento de soluções automáticas capazes de realizar tarefas complexas que antes dependiam exclusivamente da percepção e do julgamento humano. Entre essas técnicas, destacam-se os modelos de detecção de objetos em tempo real, cujo marco evolutivo é a família de algoritmos YOLO (\textit{You Only Look Once}), amplamente adotada em aplicações industriais, de vigilância e de monitoramento de ambientes \cite{redmon2016you, terven2023comprehensive}.

A verificação e a auditoria de conformidade física em ambientes de trabalho ou produção constituem etapas críticas para garantir que normas regulamentadoras, padrões de qualidade ou requisitos de segurança estejam sendo estritamente seguidos. Tradicionalmente, o processo de ``inspeção de conformidade'' consiste na verificação visual sistemática de um cenário por engenheiros, técnicos ou inspetores treinados, buscando atestar a presença de elementos regulamentares (como equipamentos de segurança, sinalização ou maquinário correto) e a ausência de irregularidades ou riscos estruturais.

No entanto, a inspeção exclusivamente manual e visual está intrinsecamente sujeita a falhas cognitivas e humanas. Fatores como a fadiga decorrente de longas jornadas de trabalho, a sobrecarga de atenção em ambientes dinâmicos de grande extensão e a subjetividade inerente à interpretação visual humana criam margens significativas para erros \cite{rahman2021context}. Em canteiros de obras de engenharia civil, por exemplo, a vastidão física das frentes de trabalho impede que os supervisores de segurança monitorem de forma contínua e em tempo real todos os colaboradores simultaneamente. Consequentemente, não conformidades como o não uso de Equipamentos de Proteção Individual (EPIs) ou a ausência de extintores de incêndio e sinalizações obrigatórias podem passar despercebidas, gerando graves riscos à integridade física do trabalhador e passivos jurídicos expressivos para as corporações.

Nesse panorama, define-se o seguinte problema de pesquisa: \textbf{Como automatizar a verificação de conformidade em ambientes por meio da detecção inteligente de objetos e análise lógica contextual a partir de imagens digitais?}

A hipótese que norteia este trabalho é que a integração de modelos de visão computacional (YOLO) com uma infraestrutura de comunicação baseada em APIs RESTful e aplicações móveis de captura permite não apenas a localização e classificação precisa de objetos em tempo real, mas também a avaliação lógica automatizada de regras de conformidade do cenário, reduzindo substancialmente a latência de inspeção e as taxas de falha humana.

Cabe salientar que, embora a segurança do trabalho em canteiros de obras de engenharia civil (como a detecção de EPIs ou de elementos de sinalização de segurança) seja adotada como o principal cenário experimental e de validação deste projeto, a solução desenvolvida foi projetada sob uma arquitetura de microsserviços modular e genérica. Isso significa que o modelo e o motor lógico de análise contextual podem ser reconfigurados e aplicados a diversos outros contextos onde a verificação de conformidade visual se faça necessária, tais como auditorias de qualidade em linhas de montagem industriais, contagem de estoque comercial, verificação de arranjo físico em salas hospitalares ou auditoria de equipamentos em TI.

O objetivo geral deste trabalho é desenvolver, implementar e testar um sistema completo de visão computacional voltado para a inspeção de conformidade de cenas. O sistema é composto por um aplicativo cliente Android escrito em Kotlin, um serviço backend FastAPI de processamento distribuído com suporte a autenticação por token JWT de duas camadas, e uma arquitetura YOLO integrada a um analisador contextual que identifica ausências ou inconsistências lógicas no cenário.

Para estruturar a validação científica deste trabalho, os objetivos específicos contemplam:
\begin{itemize}
    \item Realizar um levantamento bibliográfico acerca do estado da arte de algoritmos YOLO e sistemas de monitoramento automatizados;
    \item Elaborar e preparar conjuntos de dados personalizados para o treinamento de um modelo YOLO com foco em elementos de segurança do trabalho;
    \item Desenvolver a API RESTful backend responsável por receber dados multimídia, autenticar requisições de forma robusta e gerir a inferência;
    \item Desenvolver uma aplicação cliente móvel para a plataforma Android que capture e envie imagens/vídeos de forma performática;
    \item Analisar a performance e a precisão do sistema sob métricas estatísticas formais de avaliação visual (Precisão, Revocação e F1-score).
\end{itemize}
